\documentclass{ximera}


\begin{document}
\author{Alexander Holvoet}
\xmtitle{Zandhopen: kegels, piramides en kegelsneden}{}


\subsection*{Inleiding}

In deze lesactiviteit onderzoek je de vorm van zandhopen. Door zand (of zout, rijst, \ldots)
te gieten op verschillende basisvormen ontdek je hoe meetkunde \textemdash{} rusthoeken,
bissectrices en zelfs kegelsneden \textemdash{} opduikt in een heel concreet experiment.
Je werkt het best telkens zo: voorspel eerst wat je verwacht, voer dan het experiment uit,
en verklaar ten slotte wat je ziet.

Deze activiteit is gebaseerd op het Uitwiskeling-artikel ``Moet er nog zand zijn?''
(Uitwiskeling 42/2). De bronnen staan onderaan vermeld.

\textbf{Wat je nodig hebt:} droog zand, zout of rijst, en enkele basisvormen uit karton
(rond, vierkant, rechthoekig, driehoekig, \ldots). Plaats je basisvorm wat hoger
(bijvoorbeeld op een omgekeerd bekertje), zodat het teveel aan zand er langs de zijkant
kan afglijden.


\subsection*{Een hoop zand op een ronde ondergrond}

Giet zo veel mogelijk zand of zout op een ronde basisvorm, bijvoorbeeld een rond
bierviltje of een omgekeerd plastieken bekertje. Je moet het niet al te voorzichtig doen:
wat te veel is, glijdt er af. Daarom staat je basisvorm best een beetje hoger.

\begin{exercise}
    Welke vorm hebben de hopen?
    \begin{oplossing}
        Het zijn allemaal kegels. Je zult merken dat deze kegels niet noodzakelijk perfect
        zijn: de top is niet juist \'e\'en punt. We werken hier met een \textit{wiskundig
        model} dat vrij goed aansluit bij de realiteit, maar er niet helemaal mee samenvalt.
    \end{oplossing}
\end{exercise}

We definiëren een nieuw begrip: de \emph{rusthoek} is de hoek tussen de beschrijvende
rechten van de kegel en de horizontale richting. Die hoek is over de hele flank gelijk,
dus niet enkel vanonder.

% TODO (afbeelding ontbreekt): in de bron stond hier \afbeelding[5cm]{Lrusthoek} ter
% illustratie van de rusthoek. Deze lijntekening zit niet (meer) in img/. Voeg ze toe
% wanneer beschikbaar.

\begin{exercise}
    Denk even na hoe je de rusthoek zo nauwkeurig mogelijk kunt meten.
    \begin{hint}
        Een geodriehoek lijkt voor de hand te liggen, maar zodra je de hoop raakt, begint
        die te schuiven. Kun je de hoek meten \emph{zonder} de zandhoop aan te raken?
    \end{hint}
    \begin{oplossing}
        Een geodriehoek gebruiken is lastig: zodra je de hoop raakt, begint die te schuiven.
        Beter neem je een foto van de kegel in vooraanzicht en meet je daarop de rusthoek,
        of je plakt de foto in GeoGebra en meet de rusthoek daar. Je kunt ook de straal
        van het grondvlak en de hoogte van de kegel meten en daarmee de tangens van de
        rusthoek \textemdash{} en dus de rusthoek zelf \textemdash{} berekenen.
    \end{oplossing}
\end{exercise}

% TODO (afbeelding ontbreekt): in de bron stond hier
% \afbeelding[6cm]{rusthoek meten met geogebra}. Niet aanwezig in img/.

\begin{exercise}
    Meet en noteer de rusthoeken voor hopen met dezelfde diameter maar met verschillend
    materiaal. Verdeel eventueel het werk. Wat besluit je?
    \begin{oplossing}
        Een voorbeeld van meetresultaten:
        \begin{center}
        \begin{tabular}{l l l}
            \textbf{diameter} & \textbf{materiaal} & \textbf{rusthoek} \\
            \(10\) cm & zand  & \(32^\circ\) \\
            \(10\) cm & rijst & \(25^\circ\) \\
            \(10\) cm & zout  & \(30^\circ\) \\
        \end{tabular}
        \end{center}
        Als je de diameter van het grondvlak constant houdt, dan hangt de rusthoek af van
        het soort materiaal.
        \begin{image}[6cm]
            \includegraphics[width=6cm,keepaspectratio]{UW4202/img/lmateriaal.jpg}
        \end{image}
        Verschillende materialen geven een verschillende rusthoek.
    \end{oplossing}
\end{exercise}

\begin{exercise}
    Doe nu hetzelfde voor hopen van hetzelfde materiaal maar met verschillende diameters.
    Meet en noteer ook hier je meetresultaten in een tabel. Wat besluit je?
    \begin{oplossing}
        Een voorbeeld van meetresultaten:
        \begin{center}
        \begin{tabular}{l l l}
            \textbf{materiaal} & \textbf{diameter} & \textbf{rusthoek} \\
            zand & \(10\) cm & \(32^\circ\) \\
            zand & \(15\) cm & \(32^\circ\) \\
            zand & \(20\) cm & \(32^\circ\) \\
        \end{tabular}
        \end{center}
        De rusthoek hangt niet af van de diameter van het grondvlak.
    \end{oplossing}
\end{exercise}


\subsection*{Een hoop zand op een vierkant}

\begin{exercise}
    Wat denk je te krijgen als je zand giet op een vierkant plaatje? Doe een voorspelling
    en voer daarna pas het experiment uit.
    \begin{oplossing}
        Je krijgt een rechte piramide met het vierkant als grondvlak.
    \end{oplossing}
\end{exercise}

% TODO (afbeelding ontbreekt): in de bron stond hier \afbeelding[5cm]{LPiramide}
% (label fig:piramide). Lijntekening niet aanwezig in img/.

\begin{exercise}
    Meet opnieuw de rusthoek en vergelijk met de rusthoek bij de kegel.
    \begin{oplossing}
        De rusthoek is dezelfde als bij de kegel, tenminste als je hetzelfde materiaal
        gebruikt hebt.
    \end{oplossing}
\end{exercise}

\begin{exercise}
    Teken het bovenaanzicht van de zandhopen met een rond grondvlak en met een vierkant
    grondvlak.
    \begin{hint}
        Wat gebeurt er met de top van de kegel in bovenaanzicht? En met de ribben van de
        piramide?
    \end{hint}
    \begin{oplossing}
        Bij het ronde grondvlak krijg je een cirkel met een puntje. Bij het vierkante
        grondvlak krijg je een vierkant met diagonalen. Het puntje staat voor de top van
        de kegel, de diagonalen staan voor de ribben van de piramide. Uit het bovenaanzicht
        kun je niet afleiden hoe groot de rusthoek is.
    \end{oplossing}
\end{exercise}

% TODO (afbeelding ontbreekt): in de bron stond hier \afbeelding[7cm]{Lbovenzichtkegel}
% (label fig:bovenzicht). Lijntekening niet aanwezig in img/.


\subsection*{Een hoop zand op een rechthoekig grondvlak}

\begin{exercise}
    Welke zandhoop verwacht je bij een rechthoekig grondvlak? Maak een schets van het te
    verwachten bovenaanzicht en voer daarna pas het experiment uit.
    \begin{oplossing}
        Je krijgt geen piramide meer. Er is niet \'e\'en top maar een lijnstuk met hoogste
        punten. (In de bouw kom je dit soort ruimtelijk lichaam tegen in de vorm van een
        \textit{schilddak}.)
    \end{oplossing}
\end{exercise}

\begin{exercise}
    Wat kun je zeggen over de rusthoek?
    \begin{oplossing}
        Deze is weer constant en gelijk aan die bij de kegel en de piramide met vierkant
        grondvlak. De rusthoek hangt immers enkel af van het soort materiaal.
    \end{oplossing}
\end{exercise}

Hoogtelijnen ken je uit de aardrijkskundeles: het zijn lijnen op een kaart die alle punten
met een gelijke hoogte boven het zeeniveau verbinden.

\begin{exercise}
    Vul je tekening van het bovenaanzicht aan met enkele hoogtelijnen.
    \begin{oplossing}
        De hoogtelijnen zijn evenwijdig met de zijden van de rechthoek en de afstand tussen
        twee opeenvolgende hoogtelijnen is altijd dezelfde. Dit komt omdat de flanken overal
        even steil zijn.
    \end{oplossing}
\end{exercise}

% TODO (afbeelding ontbreekt): in de bron stond hier \afbeelding[6cm]{LRHhoogtelijnen.png}.

De hoogtelijn met de hoogste punten noemen we \textemdash{} net zoals bij echte bergen
\textemdash{} de \emph{kamlijn}. Ze is evenwijdig met de langste zijde van het grondvlak
en staat loodrecht op de kortste zijde. Je ziet dat die in het midden ligt. Het kan zijn
dat je die niet getekend hebt. Verderop proberen we de positie van de kamlijn meetkundig
te bepalen.

\begin{exercise}
    Verbind de hoekpunten van de hoogtelijnen.
    \begin{oplossing}
        Je krijgt een figuur waarin door elk hoekpunt van de hoogtelijnen een schuine
        rechte loopt onder \(45^\circ\) met de zijden.
    \end{oplossing}
\end{exercise}

% TODO (afbeelding ontbreekt): in de bron stond hier \afbeelding[6cm]{LRHhoogtelijnenbiss.png}.

\begin{exercise}
    Aan welke eigenschap voldoen
    \begin{enumerate}[label=(\alph*)]
        \item punten die op de kamlijn liggen, zoals het punt \(R\);
        \item punten op \'e\'en van de schuine rechten die de hoekpunten van de
            rechthoekige hoogtelijnen verbinden, zoals het punt \(P\);
        \item punten die op zowel de kamlijn als op de twee schuine rechten liggen, zoals
            het punt \(S\)?
    \end{enumerate}
    \begin{hint}
        Denk telkens aan de afstand van het punt tot de randen van het grondvlak. Hoe ver
        ligt een zandkorrel op dit punt van de verschillende randen?
    \end{hint}
    \begin{oplossing}
        \begin{enumerate}[label=(\alph*)]
            \item Punten op de kamlijn bevinden zich even ver van de twee overstaande
                zijden, zowel in 3D als in het bovenaanzicht. Dit verklaart waarom de
                kamlijn in het bovenaanzicht in het midden tussen de twee lange randen ligt.
            \item Deze punten liggen even ver van twee opeenvolgende randen, ook weer zowel
                in 3D als in bovenaanzicht. Dit verklaart waarom deze lijnen een hoek van
                \(45^\circ\) maken met deze zijden.
            \item Deze punten bevinden zich even ver van drie opeenvolgende randen.
        \end{enumerate}
    \end{oplossing}
\end{exercise}

% TODO (afbeelding ontbreekt): in de bron stond hier \afbeelding[6cm]{LRHbiss.png}
% (met de punten R, P en S erop aangeduid).

In het bovenaanzicht is de \textit{kamlijn} dus de meetkundige plaats van de punten die
zich even ver bevinden van de twee dichtstbijzijnde zijden. De punten op het uiteinde van
de kamlijn liggen op gelijke afstand van drie zijden. In 3D kunnen we de ribben daarom ook
kamlijnen noemen.


\subsection*{Een hoop zand op een driehoekig grondvlak}

\begin{exercise}
    Neem als volgend grondvlak een willekeurige driehoek. Schets eerst wat je verwacht als
    bovenaanzicht, en teken ook de vermoedelijke hoogtelijnen. Giet dan zoveel mogelijk
    zand op de omhoog geplaatste driehoek. Komt het experiment overeen met je schets? Is de
    rusthoek ook hier weer gelijk aan die in de vorige experimenten?
    \begin{oplossing}
        Je krijgt een piramide, zoals je in de foto's hieronder ziet. Die vorm hangt enkel
        af van het grondvlak en niet van waar je begint met zand te gieten. De rusthoek is
        weer dezelfde.
        \begin{image}[12cm]
            \includegraphics[width=6cm,keepaspectratio]{UW4202/img/lzandopdriehoek.jpg}
            \quad
            \includegraphics[width=6cm,keepaspectratio]{UW4202/img/lsituatie2.jpg}
        \end{image}
        Een zandhoop op een driehoekig grondvlak: een piramide met de hoogtelijnen erop.
    \end{oplossing}
\end{exercise}

% Opmerking voor de leerkracht: in de bron stond hiernaast ook een lijntekening van de
% hoogtelijnen op de driehoek (Lhoogtelijnendriehoek.png), niet aanwezig in img/.

\begin{exercise}
    Als je ook hier de hoekpunten van de hoogtelijnen verbindt, zie je speciale lijnen
    verschijnen op het bovenaanzicht. Welke lijnen zijn dit?
    \begin{oplossing}
        Het zijn de bissectrices van de hoeken gevormd door twee opeenvolgende zijden van
        de driehoek.
    \end{oplossing}
\end{exercise}

% TODO (afbeelding ontbreekt): \afbeelding[6cm]{Ldriehoekkortsteweg.png}.

\begin{exercise}
    Geef hiervoor een verklaring. Waarmee komen deze lijnen overeen in 3D?
    \begin{oplossing}
        De bissectrices zijn de meetkundige plaats van de punten die even ver liggen van
        twee opeenvolgende zijden van de driehoek. In 3D komen ze overeen met de ribben van
        de piramide. Korrels op de ribben weten niet naar welke kant ze moeten vallen en
        blijven dus liggen.
    \end{oplossing}
\end{exercise}

\begin{exercise}
    We weten dat de drie bissectrices van een driehoek elkaar in \'e\'en punt snijden. Wat
    weet je van dit snijpunt?
    \begin{oplossing}
        Omdat het snijpunt op elk van de bissectrices ligt, ligt het even ver van de drie
        zijden van de driehoek. De loodlijnen op de zijden geven de kortste weg aan voor een
        zandkorrel om naar beneden te rollen.
    \end{oplossing}
\end{exercise}

% TODO (afbeelding ontbreekt): \afbeelding[8cm]{zandopdriehoek.png}.

\begin{exercise}
    Wat betekent het snijpunt van de bissectrices uit het bovenaanzicht voor de zandhoop?
    \begin{oplossing}
        Dit snijpunt is de projectie van de top van de zandhoop. Deze top is het punt dat
        zich even ver van de drie randen van het grondvlak bevindt.
    \end{oplossing}
\end{exercise}


\subsection*{Convexe grondvlakken: trapezia}

\begin{instructorNotes}
    Dit stuk is in de bron niet volledig als lesactiviteit uitgeschreven, maar het is wel
    de bedoeling om dit op een onderzoekende manier aan te brengen, zoals in de vorige
    paragrafen. Dit kan in groepjes of via een onderwijsleergesprek, omdat het wat
    moeilijker wordt. Het interessante van deze situatie is niet zozeer het eindresultaat,
    maar vooral de zoektocht ernaartoe.
\end{instructorNotes}

We bekijken enkele trapezia. Een trapezium is een vierhoek met twee evenwijdige zijden.
We vertrekken van een driehoek: teken de bissectrices en de ingeschreven cirkel. Doe dit
in viervoud. Knip nu van drie van die driehoeken een stuk af volgens een lijn evenwijdig
met de basis: boven de ingeschreven cirkel (situatie 2), rakend aan de ingeschreven cirkel
(situatie 3) en dwars door de ingeschreven cirkel (situatie 4). Teken ten slotte ook de
bissectrices van de twee nieuw verschenen hoeken van elk trapezium.

% TODO (afbeelding ontbreekt): \afbeelding[15cm]{Ltrapezia.png} (een driehoek en drie
% soorten trapezia). Lijntekening niet aanwezig in img/.

Voor je het experiment uitvoert, kun je proberen te voorspellen welke zandhoop en welk
bovenaanzicht je zult krijgen. Het loont echt de moeite om dit zelf verder uit te zoeken.
Als je zand giet op de volledige driehoek (situatie 1), krijg je uiteraard een piramide.
We bekijken situatie 2 en 4 in detail; situatie 3 laten we aan de lezer.

\textbf{Situatie 2.} In de tweede situatie zijn de bissectrices niet concurrent. Er zijn
twee snijpunten van telkens twee bissectrices, de punten \(O\) en \(O'\). Ze liggen op de
bissectrice van de afgesneden hoek van de driehoek en zijn de projecties van \(S\) en
\(S'\), de eindpunten van de kamlijn.

In de 3D-tekening zien we twee rechthoekige driehoeken \(HSO\) en \(H'S'O'\) die dezelfde
rusthoek hebben en dus gelijkvormig zijn. We zien dat \(|OS| = |OH|\cdot \tan\theta\) en
\(|O'S'| = |O'H'|\cdot \tan\theta\), maar \(|OH| > |O'H'|\), dus \(|OS| > |O'S'|\). Het
lijnstuk \([SS']\) loopt dus niet horizontaal. In het bovenaanzicht zie je dit aan het feit
dat \(O'\) dichter bij de rand ligt dan \(O\). Verifieer het met het experiment. Sterker
nog:
\[
    |OS| - |O'S'| = (|OH| - |O'H'|)\cdot \tan\theta .
\]
Als we de kniplijn helemaal zouden opschuiven tot in het hoekpunt, dan zou \(|O'H'| = 0\)
worden, dus de helling van \([SS']\) zou dezelfde zijn als die van \([SA]\). Als we
\([SS']\) verlengen, zou die door de top \(A\) van de afgeknipte driehoek lopen.

\begin{image}[7cm]
    \includegraphics[width=7cm,keepaspectratio]{UW4202/img/lsituatie2.jpg}
\end{image}
Situatie 2: het afgeknipte trapezium met de kamlijn.

% TODO (afbeeldingen ontbreken): in de bron stonden hier de lijntekeningen
% Ltrap2.png (projectie situatie 2) en Ltrap23D.png (3D-versie). Niet aanwezig in img/.

\textbf{Situatie 4.} Hier hebben we ook de vier bissectrices van de hoeken van het
trapezium getekend. Nu snijdt telkens \'e\'en van de oude bissectrices met \'e\'en van de
nieuwe bissectrices. Dit geeft weer twee punten \(O\) en \(O'\). Het lijnstuk \([OO']\) dat
deze punten verbindt, is evenwijdig met de twee evenwijdige zijden van het trapezium.
\(O\) en \(O'\) zijn de projecties op het grondvlak van de toppen \(S\) en \(S'\).

In de 3D-voorstelling zijn er weer twee driehoeken \(HSO\) en \(H'S'O'\) die dezelfde
rusthoek met het grondvlak hebben. We vinden \(|OS| = |OH|\cdot \tan\theta\) en
\(|O'S'| = |O'H'|\cdot \tan\theta\), maar nu is \(|OH| = |O'H'|\) en dus \(|OS| = |O'S'|\).
Bijgevolg is \([SS'] \parallel [HH'] \parallel [KK']\): de kamlijn is evenwijdig met de
evenwijdige zijden van het trapezium. Verifieer weer met het experiment.

\begin{image}[7cm]
    \includegraphics[width=7cm,keepaspectratio]{UW4202/img/lsituatie4.jpg}
\end{image}
Situatie 4: de kamlijn is evenwijdig met de evenwijdige zijden van het trapezium.

% TODO (afbeeldingen ontbreken): in de bron stonden hier Ltrap4.png (projectie situatie 4)
% en Ltrap43D.png (3D-tekening). Niet aanwezig in img/.


\subsection*{Een oefening op convexe veelhoeken}

\begin{exercise}
    Construeer, door gebruik te maken van bissectrices, de bovenaanzichten van de zandhopen
    met de gegeven grondvlakken. Je kunt hiervoor GeoGebra gebruiken. Verifieer daarna met
    een experiment.
    \begin{hint}
        De kamlijn bestaat uit de punten die even ver liggen van de twee dichtstbijzijnde
        randen. Dat zijn net stukken bissectrice tussen opeenvolgende zijden.
    \end{hint}
    \begin{oplossing}
        Teken bij elke veelhoek de bissectrices van de hoeken; de kamlijnen vind je waar de
        bissectrices van opeenvolgende zijden samenkomen. De aardrijkskundige hoogtelijnen
        helpen om de constructie te controleren, maar zijn niet nodig om de oefening op te
        lossen.
    \end{oplossing}
\end{exercise}

% TODO (afbeeldingen ontbreken): in de bron stonden hier Lconvexeveelhoeken.png (de op te
% lossen grondvlakken) en Lconvexemethoogtel.png (de oplossingen met hoogtelijnen). Niet
% aanwezig in img/.


\subsection*{Zand op een schijf met een gat in}

Het grondvlak waarop we zand willen gieten, is een cirkelschijf met een gaatje in. Je ziet
hieronder twee cirkelschijven: \'e\'en met een gat in het midden, \'e\'en met een
gedecentreerd gat. Het gat is klein en in de redeneringen behandelen we het als een punt.
In werkelijkheid mag het gat ook niet \emph{te} klein zijn, want de zandkorrels moeten
erdoor kunnen.

\begin{image}[12cm]
    \includegraphics[width=5.5cm,keepaspectratio]{UW4202/img/zand07.jpg}
    \quad
    \includegraphics[width=5.5cm,keepaspectratio]{UW4202/img/zand08.jpg}
\end{image}
Links een centraal gat, rechts een gedecentreerd gat.

\begin{exercise}
    Teken op deze bovenaanzichten telkens de kamlijn, giet dan zand op het kartonnen
    exemplaar om dit te verifi\"eren, en verklaar ten slotte de vorm van de kamlijn.
    Begin met het \textbf{centrale} gat.
    \begin{hint}
        De kamlijn bestaat uit de punten die even ver liggen van het gat als van de rand
        van de schijf. Welke kromme krijg je bij een centraal gat?
    \end{hint}
    \begin{oplossing}
        Bij het centrale gat is het bovenaanzicht van de kamlijn een concentrische cirkel
        met als straal de helft van de straal van de gegeven schijf. Deze cirkel bestaat
        immers uit de punten op gelijke afstand tot het gat en de rand. In 3D heb je in het
        midden van de kegel een kegelvormige uitsparing met dezelfde rusthoek: de zandhoop
        is een kegel waaruit een tweede, omgekeerde kegel met dezelfde rusthoek is
        weggehaald. Die tweede kegel bepaalt de uitsparing.
        \begin{image}[12cm]
            \includegraphics[width=4cm,keepaspectratio]{UW4202/img/zand09.jpg}
            \quad
            \includegraphics[width=6cm,keepaspectratio]{UW4202/img/zand10.jpg}
        \end{image}
        Links het boven- en vooraanzicht, rechts de kegel met de omgekeerde kegel die de
        uitsparing bepaalt.
    \end{oplossing}
\end{exercise}

\begin{exercise}
    Bekijk nu het \textbf{gedecentreerde} gat. Welke kromme is de kamlijn hier in
    bovenaanzicht? Bewijs je vermoeden.
    \begin{hint}
        Een punt \(P\) van de kamlijn ligt even ver van het gat \(G\) als van de rand van
        de schijf. Meet de afstand tot de rand loodrecht, dus langs een straal. Wat kun je
        zeggen over \(|PM| + |PG|\), met \(M\) het middelpunt van de schijf?
    \end{hint}
    \begin{oplossing}
        De kamlijn is in bovenaanzicht een \emph{ellips}. Als het gat dicht bij het
        middelpunt ligt, lijkt het op het eerste gezicht een cirkel, maar laten we
        redeneren. Een willekeurig punt \(P\) van de kamlijn ligt even ver van het gat \(G\)
        als van de rand van de schijf. Die afstand \(|PQ|\) tot de rand wordt loodrecht op
        die rand gemeten, dus langs een straal. Bijgevolg is
        \[
            |PM| + |PG| = |PM| + |PQ| = |MQ| = r ,
        \]
        waarbij \(r\) de straal van de cirkel is. De som \(|PM| + |PG|\) is dus constant, en
        de punten \(P\) van de kamlijn vormen een ellips met als brandpunten \(G\) en \(M\).
        \begin{image}[4cm]
            \includegraphics[width=4cm,keepaspectratio]{UW4202/img/zand11.jpg}
        \end{image}
    \end{oplossing}
\end{exercise}

\begin{exercise}
    Je kunt, bij het gedecentreerde gat, ook redeneren met hoogtelijnen. Je ziet hier de
    cirkelschijf met het gat, en concentrische cirkels rond het middelpunt en rond het gat.
    Construeer met deze figuur de kamlijn en leg uit.
    \begin{image}[4cm]
        \includegraphics[width=4cm,keepaspectratio]{UW4202/img/zand12.jpg}
    \end{image}
    \begin{oplossing}
        Je kunt in het vlak redeneren: de punten van de kamlijn liggen op gelijke afstand
        van het gat als van de rand. Bijvoorbeeld 5 eenheden van het gat en 5 eenheden van
        de rand, waarbij een eenheid de afstand tussen twee concentrische cirkels is. Dit
        geeft de punten \(P\) en \(P'\). Analoog voor andere aantallen eenheden.

        Je kunt ook in 3D redeneren: de cirkels zijn hoogtelijnen, en je snijdt de cirkels
        die even hoog zijn. De cirkels rond het gat beginnen in het gat met hoogte nul
        (zoals de kegel op zijn top); de cirkels rond het middelpunt beginnen aan de rand
        met hoogte nul (zoals de gewone kegel).
        \begin{image}[4cm]
            \includegraphics[width=4cm,keepaspectratio]{UW4202/img/zand13.jpg}
        \end{image}
    \end{oplossing}
\end{exercise}

\begin{exercise}
    Het gat was tot nu toe heel klein en we behandelden het als een punt. Wat verandert er
    als het (gedecentreerde) gat een grotere cirkel(schijf) is?
    \begin{oplossing}
        Het bovenaanzicht van de kamlijn is opnieuw een ellips. De brandpunten zijn het
        middelpunt \(M\) van de grote cirkel en het middelpunt \(G\) van het gat. De som van
        de afstanden tot deze twee punten is weer constant, maar deze constante is nu niet
        de straal van de cirkel, wel de som van de straal van de cirkel en die van het gat.

        Een punt \(P\) van de kamlijn ligt even ver van de rand van de grote cirkel als van
        de rand van het ronde gat: \(|PQ| = |PH|\). De som van de afstanden van \(P\) tot
        \(M\) en \(G\) is bijgevolg
        \[
            |PM| + |PG| = |PM| + |PH| + |HG| = |PM| + |PQ| + |HG| = |MQ| + |HG| = r + s ,
        \]
        waarbij \(r\) de straal van de grote cirkel is en \(s\) de straal van het gat. Deze
        som is constant, dus \(P\) beschrijft een ellips.
        \begin{image}[12cm]
            \includegraphics[width=4cm,keepaspectratio]{UW4202/img/zand14.jpg}
            \quad
            \includegraphics[width=6cm,keepaspectratio]{UW4202/img/zand15.jpg}
        \end{image}
    \end{oplossing}
\end{exercise}


\subsection*{Zand op een cirkelsegment}

\begin{image}[5cm]
    \includegraphics[width=5cm,keepaspectratio]{UW4202/img/zand16.jpg}
\end{image}
Een cirkelsegment als grondvlak.

We nemen als basis voor een zandhoop een deel van een cirkelschijf aan \'e\'en kant van een
koorde. Deze figuur noemen we een \emph{cirkelsegment}. Het weggesneden stuk is ook een
cirkelsegment. We nemen hier een `groot' cirkelsegment \textemdash{} dit wil zeggen het
stuk met het middelpunt erop \textemdash{} maar het gaat even goed met een `klein'
cirkelsegment.

\begin{exercise}
    Teken (het bovenaanzicht van) de kamlijn die je verwacht, giet dan zand op het
    kartonnen exemplaar om dit te verifi\"eren, en bepaal waaraan de punten van de kamlijn
    voldoen.
    \begin{oplossing}
        Voor elk punt \(P\) van de kamlijn is de afstand \(|PQ|\) tot de koorde gelijk aan
        de afstand \(|PR|\) tot de rand van de cirkel.
        \begin{image}[5cm]
            \includegraphics[width=5cm,keepaspectratio]{UW4202/img/zand17.jpg}
        \end{image}
    \end{oplossing}
\end{exercise}

Je vermoedt waarschijnlijk dat (het bovenaanzicht van) de kamlijn een parabool is. In de
volgende opgaven bewijs je dit.

\begin{exercise}
    Om te bewijzen dat de kamlijn een parabool is, heb je een brandpunt en een richtlijn
    nodig. Het brandpunt is het middelpunt \(M\) van de cirkel. Waar ongeveer moet de
    richtlijn \(d\) liggen? Toon dit met je vinger.
    \begin{oplossing}
        De richtlijn is een rechte evenwijdig met de koorde (vanwege de symmetrie) en
        `onder' de kamlijn. De precieze plaats komt in de volgende vraag aan bod.
    \end{oplossing}
\end{exercise}

\begin{exercise}
    Bewijs, door een geschikte keuze van de richtlijn, dat de kamlijn inderdaad een
    parabool is.
    \begin{hint}
        Neem een willekeurig punt \(P\) van de kamlijn en noem \(r\) de straal. Gebruik dat
        \(|PR| = |PQ|\), dus \(|PM| = r - |PR| = r - |PQ|\). Op welke afstand van de koorde
        moet je \(d\) leggen opdat \(|PM| = |PS|\)?
    \end{hint}
    \begin{oplossing}
        Neem een willekeurig punt \(P\) van de kamlijn en noem \(r\) de straal van de
        cirkel. We bepalen de richtlijn \(d\) zo dat de afstand \(|PM|\) tot het brandpunt
        gelijk is aan de afstand \(|PS|\) tot de rechte \(d\). We weten dat \(|PR| = |PQ|\),
        en dus
        \[
            |PM| = r - |PR| = r - |PQ| .
        \]
        Om \(|PM| = |PS|\) te krijgen, moeten we \(d\) op afstand \(r\) van de koorde nemen.
        Dan geldt inderdaad
        \[
            |PM| = r - |PR| = r - |PQ| = |QS| - |PQ| = |PS| .
        \]
        \begin{image}[5cm]
            \includegraphics[width=5cm,keepaspectratio]{UW4202/img/zand18.jpg}
        \end{image}
    \end{oplossing}
\end{exercise}

\begin{exercise}
    Construeer de kamlijn met hoogtelijnen.
    \begin{image}[5cm]
        \includegraphics[width=5cm,keepaspectratio]{UW4202/img/zand19.jpg}
    \end{image}
    \begin{oplossing}
        Construeer de snijpunten van de overeenkomstige hoogtelijnen. De rechte hoogtelijnen
        horen bij een schuin vlak dat volgens de rusthoek afloopt naar de koorde van het
        grondvlak. De cirkelvormige hoogtelijnen horen bij de kegel die je zou hebben als de
        basis de volledige cirkelschijf was. De kamlijn is de doorsnede van die kegel met
        het schuine vlak (even schuin als de beschrijvenden van de kegel). Dit is nog een
        extra manier om in te zien dat het een stuk parabool is.
        \begin{image}[15cm]
            \includegraphics[width=5cm,keepaspectratio]{UW4202/img/zand20.jpg}
            \quad
            \includegraphics[width=5cm,keepaspectratio]{UW4202/img/zand21.jpg}
            \quad
            \includegraphics[width=5cm,keepaspectratio]{UW4202/img/zand22.jpg}
        \end{image}
    \end{oplossing}
\end{exercise}


\subsection*{Zand op een niet-convexe veelhoek}

\begin{exercise}
    Teken de kamlijnen die je verwacht als je op deze figuur zand zou gieten.
    \begin{image}[7cm]
        \includegraphics[width=7cm,keepaspectratio]{UW4202/img/zand26.jpg}
    \end{image}
    \begin{hint}
        Je eerste poging bestaat wellicht uit rechte lijnstukken (bissectrices). Verifieer
        met zand: in de binnenbocht klopt er iets niet. De zandhoop komt daar hoger dan je
        verwacht.
    \end{hint}
    \begin{oplossing}
        Veel leerlingen tekenen eerst enkel rechte lijnstukken.
        \begin{image}[12cm]
            \includegraphics[width=6cm,keepaspectratio]{UW4202/img/zand27.jpg}
            \quad
            \includegraphics[width=6cm,keepaspectratio]{UW4202/img/zand30.jpg}
        \end{image}
        Wie dit met zand verifieert, ziet dat de kamlijn niet alleen uit rechte lijnstukken
        bestaat: in zijaanzicht komt de zandhoop in de bocht wat hoger.
        \begin{image}[12cm]
            \includegraphics[width=5cm,keepaspectratio]{UW4202/img/zand28.jpg}
            \quad
            \includegraphics[width=5cm,keepaspectratio]{UW4202/img/zand29.jpg}
        \end{image}
        De stukken van de kamlijn kunnen niet allemaal rechtlijnig zijn. De punten tussen
        \(K\) en \(M\) (\(M\) zelf niet meegerekend) liggen dichter bij rand \(d\) dan bij
        punt \(F\). Het lijnstuk \([KM]\) moet vervangen worden door de punten die even ver
        van \(F\) als van rand \(d\) liggen: een stuk \emph{parabool} met brandpunt \(F\) en
        richtlijn \(d\). Analoog wordt \([KL]\) een stuk parabool met brandpunt \(F\) en
        richtlijn \(e\). Het lijnstuk \([KF]\) moet weggelaten worden: dit is geen kamlijn.
        Zandkorrels boven \([KF]\) twijfelen niet en rollen naar \(F\), het dichtste
        randpunt.
        \begin{image}[12cm]
            \includegraphics[width=6cm,keepaspectratio]{UW4202/img/zand31.jpg}
            \quad
            \includegraphics[width=6cm,keepaspectratio]{UW4202/img/zand32.jpg}
        \end{image}
        Met hoogtelijnen zie je dat de zandhoop bij het snijpunt van de twee paraboolstukken
        inderdaad iets hoger ligt dan de rechte kamlijnen.
    \end{oplossing}
\end{exercise}


\subsection*{Zand op een ellips}

\begin{exercise}
    Giet zoveel mogelijk droog zand op een ellips die in karton is uitgeknipt en horizontaal
    op een bekertje rust. Wat zie je?
    \begin{oplossing}
        In bovenaanzicht zien we een kamlijn: een stuk van de grote as, symmetrisch rond het
        middelpunt. Deze kamlijn lijkt niet tot aan de toppen op de grote as te gaan. Een
        eerste hypothese: misschien vallen de eindpunten van (het bovenaanzicht van) de
        kamlijn samen met de brandpunten van de ellips? In vooraanzicht zien we dat deze kam
        niet horizontaal loopt, maar in het midden hoger is dan opzij. Dat is niet
        verwonderlijk: de punten in het midden van de ellips liggen verder van de rand.
        \begin{image}[12cm]
            \includegraphics[width=6cm,keepaspectratio]{UW4202/img/zand35.jpg}
            \quad
            \includegraphics[width=6cm,keepaspectratio]{UW4202/img/zand36.jpg}
        \end{image}
    \end{oplossing}
\end{exercise}

Vanaf nu werken we in het bovenaanzicht. De punten van de kamlijn beschouwen we als punten
van de grote as van de ellips. We kunnen verwijzen naar de afrollende zandkorrels in de
ruimte, maar als we punten of rechten benoemen, doen we aan vlakke meetkunde.

\begin{exercise}
    Net als in de vorige opgaven kun je de kamlijn bekijken als een meetkundige plaats van
    punten die aan een bepaalde voorwaarde voldoen. Aan welke voorwaarde moeten de punten
    van de kamlijn voldoen?
    \begin{oplossing}
        De kamlijn is de meetkundige plaats van de punten op gelijke loodrechte afstand tot
        twee of meer randen van de figuur, hier dus tot de ellips. Zandkorrels rollen immers
        naar de dichtste rand; korrels aan de ene kant van de kamlijn rollen naar de ene
        rand, korrels aan de andere kant naar de andere rand.

        Voor de punten van de grote as \([DE]\) die niet te ver van het middelpunt liggen,
        klopt dit duidelijk. Punten van de grote as die dicht bij een top \(D\) of \(E\)
        liggen, rollen recht naar die top: zij liggen niet op een kamlijn, want er is maar
        \'e\'en dichtste randpunt. De vraag is dus: waar liggen de eindpunten \(A\) en \(B\)
        van de kamlijn, zo dat een punt \(P\) tussen \(A\) en \(B\) twee dichtste randpunten
        \(M\) en \(M'\) heeft, en een punt tussen \(A\) en \(D\) (of tussen \(B\) en \(E\))
        slechts \'e\'en, namelijk \(D\) (of \(E\))?
        \begin{image}[6cm]
            \includegraphics[width=6cm,keepaspectratio]{UW4202/img/zand37.jpg}
        \end{image}
    \end{oplossing}
\end{exercise}

\begin{exercise}
    Je kunt het misschien omkeren. In plaats van een variabel punt \(P\) op de grote as te
    nemen en je af te vragen hoe dicht het bij de top moet liggen, kun je een variabel punt
    \(M\) van de ellips nemen, dat steeds dichter naar de top \(D\) laten gaan, en bepalen
    waar het bijbehorende punt \(P\) van de grote as in de limiet naartoe gaat. Onderzoek
    dit met GeoGebra voor de ellips \(\frac{x^2}{25} + \frac{y^2}{9} = 1\).
    \begin{oplossing}
        Teken deze ellips met GeoGebra, neem een punt \(M\) erop, en gebruik de knoppen
        `raaklijn' en `loodlijn' om de normaal in \(M\) te tekenen en met de grote as te
        snijden. Voor \(M\) heel dicht bij de top \(D\) heeft het punt \(P\) coördinaat
        \((3{,}2;\, 0)\). Dit is nog geen exacte methode, maar het lijkt er al op dat het
        eindpunt \(A\) van de kam niet het brandpunt is. Deze ellips heeft brandpunten
        \(F_{1,2}(\pm 4,\, 0)\).
        \begin{image}[12cm]
            \includegraphics[width=6cm,keepaspectratio]{UW4202/img/zand38.jpg}
            \quad
            \includegraphics[width=6cm,keepaspectratio]{UW4202/img/zand39.jpg}
        \end{image}
    \end{oplossing}
\end{exercise}

\begin{exercise}
    We willen nu het eindpunt \(A\) aan de rechterkant van de kam exact lokaliseren. Ken je
    een eigenschap van de ellips waar de normaal een rol in speelt?
    \begin{oplossing}
        De weerkaatsingseigenschap zegt dat \(F_1 M\) en \(F_2 M\) gelijke hoeken vormen met
        de normaal \(MP\) (waardoor een lichtstraal die in \(F_1\) vertrekt, in \(F_2\)
        toekomt en omgekeerd). Dus: \(MP\) is de bissectrice van de hoek
        \(\widehat{F_1 M F_2}\).
        \begin{image}[8.5cm]
            \includegraphics[width=8.5cm,keepaspectratio]{UW4202/img/zand40.jpg}
        \end{image}
    \end{oplossing}
\end{exercise}

\begin{exercise}
    Wat zegt dit over de positie van punt \(P\) op de grote as? Preciezer: wat weet je over
    de verhouding \(\frac{|F_1 P|}{|F_2 P|}\)?
    \begin{hint}
        Pas de bissectricestelling toe in de driehoek \(F_1 M F_2\), met \(MP\) als
        bissectrice vanuit \(M\).
    \end{hint}
    \begin{oplossing}
        De bissectricestelling zegt: de bissectrice van een hoek van een driehoek verdeelt
        de overstaande zijde in stukken die in dezelfde verhouding staan als de aanliggende
        zijden. In de driehoek \(F_1 M F_2\) geeft dit
        \[
            \frac{|F_1 P|}{|F_2 P|} = \frac{|F_1 M|}{|F_2 M|} .
        \]
    \end{oplossing}
\end{exercise}

\begin{exercise}
    Maak de redenering af: waar komt \(P\) terecht als \(M\) naar de top \(D\) gaat? Maak
    het ten slotte concreet voor de ellips \(\frac{x^2}{25} + \frac{y^2}{9} = 1\).
    \begin{oplossing}
        Als \(M\) naar \(D\) gaat, komt \(P\) in het grenspunt \(A\) van de kamlijn
        \([AB]\). Het punt \(A\) is dus bepaald door
        \[
            \frac{|F_1 A|}{|F_2 A|} = \frac{|F_1 D|}{|F_2 D|} .
        \]
        Voor de ellips \(\frac{x^2}{25} + \frac{y^2}{9} = 1\) is de halve grote as \(5\) en
        de halve kleine as \(3\). We hebben \(D(5, 0)\), \(F_1(4, 0)\) en \(F_2(-4, 0)\).
        Stel \(A(u, 0)\); we zoeken \(u\) met
        \[
            \frac{4 - u}{u - (-4)} = \frac{5 - 4}{5 - (-4)} .
        \]
        Oplossen naar \(u\) geeft \(u = 3{,}2\), precies wat we met GeoGebra vonden. (We
        gingen ervan uit dat \(u < 4\), wat het GeoGebra-onderzoek overtuigend liet zien.)
    \end{oplossing}
\end{exercise}


\subsection*{Bronnen}

% Open beslissing (zie CONVERSIE.md): bronnen voorlopig hier opgenomen als gewone lijst.
% Te bespreken met Alexander hoe we bronnen structureel integreren in de Ximera-versie.
\begin{itemize}
    \item Dendoncker, J.-M. (2025). Moet er nog zand zijn? \emph{Wiskunde en Onderwijs, 51/203}, 28--34.
    \item Élèves des collèges Vieux Colombier (Le Mans) et Vieux Chêne. (2009). La géométrie du tas de sable. \url{http://soleil.rene.free.fr/neo/Butinages/La_geometrie_du_tas_de_sable.pdf}
    \item Eggermont, H., Roelens, M. (2004). Kegelsneden. \emph{Uitwiskeling 20/4}, 9--42.
    \item Gilbert, T. (2017). Tas de sable sur une ellipse. In Gilbert, T., Ninove, L., le GEM, \emph{Le plaisir de chercher en mathématiques}. Louvain-la-Neuve: Presses Universitaires de Louvain.
    \item Levrie, P., Penne, R. (2015). Kegelsneden made easy. \emph{Uitwiskeling 31/3}, 15--17.
    \item Matemateca. (z.d.). Sandpiles. \url{https://matemateca.ime.usp.br/acervo/montanhas_areia_ingles.html}
    \item M3 Matemática Multimídia. (23 maart 2012). Montanhas geométricas. YouTube: \url{https://www.youtube.com/watch?v=f2kBia-jNfg}
    \item Pierson, P. (2012). La géométrie des tas de sable. \emph{Travail de fin d'études}, Haute École Galilée.
\end{itemize}

\end{document}
